---
tipo: Discursiva
dificuldade: Média
nivel: médio
disciplina: Matemática
assunto: Geometria Analítica
gabarito:
resposta: (x-a)^2 + (y-a)^2 = a^2,\ \text{onde}\ a = \frac{2-\sqrt{2}}{2+\sqrt{2}}
tags: [circunferência, tangência, figura]
fonte: concurso
concurso:
banca: UFF
ano: 1999
numero: 1
cargo:
prova:
---
\question
A circunferência $C_1$, de raio 1, é tangente aos eixos coordenados, conforme representação abaixo.

\includegraphics{figura-1}

Determine a equação da circunferência $C_2$, tangente simultaneamente aos eixos coordenados e à $C_1$.

---
tipo: Discursiva
dificuldade: Média
nivel: médio
disciplina: Matemática
assunto: Trigonometria
gabarito:
resposta: ab = 1
tags: [identidades trigonométricas]
fonte: concurso
concurso:
banca: UFF
ano: 1999
numero: 2
cargo:
prova:
---
\question
Determine a relação entre os números reais $a$ e $b$ de modo que as igualdades

$1 + \cos x = a \sen x$ e $1 - \cos x = b \sen x$,

com $x \ne k\pi$, $k \in \mathbb{Z}$, sejam satisfeitas simultaneamente.

---
tipo: Discursiva
dificuldade: Média
nivel: médio
disciplina: Matemática
assunto: Teoria dos Conjuntos
gabarito:
resposta: a)\ P = \{3, 4, 5, 7\},\ Q = \{1, 2, 3, 7\},\ R = \{2, 5, 6, 7\};\ b)\ \{3\};\ c)\ \{2, 5, 7\};\ d)\ \{2, 6\};\ e)\ \{2, 3, 4, 5, 7\}
tags: [enumeração, operações entre conjuntos, figura]
fonte: concurso
concurso:
banca: UFF
ano: 1999
numero: 3
cargo:
prova:
---
\question
Considere os conjuntos representados abaixo.

\includegraphics{figura-2}

Represente, enumerando seus elementos, os conjuntos:
\begin{parts}
\part $P$, $Q$ e $R$
\part $(P \cap Q) - R$
\part $(P \cup Q) \cap R$
\part $(P \cup R) - P$
\part $(Q \cap R) \cup P$
\end{parts}

---
tipo: Discursiva
dificuldade: Fácil
nivel: médio
disciplina: Matemática
assunto: Função Afim
gabarito:
resposta: n = 45
tags: [equação do primeiro grau, lucro]
fonte: concurso
concurso:
banca: UFF
ano: 1999
numero: 4
cargo:
prova:
---
\question
Um baleiro vende $n$ balas, por R$ 0,30 cada, e obtém $L$ reais. Se vender 15 balas a menos, por R$ 0,45 cada, obterá os mesmos $L$ reais.

Determine o valor de $n$.

---
tipo: Discursiva
dificuldade: Fácil
nivel: superior
disciplina: Matemática
assunto: Polinômios
gabarito:
resposta: p(1) = 5
tags: [resto da divisão, teorema do resto]
fonte: concurso
concurso:
banca: UFF
ano: 1999
numero: 5
cargo:
prova:
---
\question
O resto da divisão do polinômio $p(x)$ por $(x-1)^3$ é o polinômio $r(x)$.

Sabendo que o resto da divisão de $r(x)$ por $x - 1$ é igual a 5, encontre o valor de $p(1)$.

---
tipo: Discursiva
dificuldade: Média
nivel: médio
disciplina: Matemática
assunto: Geometria Plana
gabarito:
resposta: \text{Área} = 100\sqrt{3}\ \text{cm}^2
tags: [retângulo, circunferência, figura]
fonte: concurso
concurso:
banca: UFF
ano: 1999
numero: 6
cargo:
prova:
---
\question
Determine a área do retângulo $MNPQ$, representado na figura abaixo, sabendo que a diagonal $MP$ é o diâmetro da circunferência $C$ cujo raio mede 10 cm.

\includegraphics{figura-3}

---
tipo: Discursiva
dificuldade: Fácil
nivel: médio
disciplina: Matemática
assunto: Progressão Aritmética
gabarito:
resposta: \text{O 3º termo negativo é } -33
tags: [sequência]
fonte: concurso
concurso:
banca: UFF
ano: 1999
numero: 7
cargo:
prova:
---
\question
Determine o terceiro termo negativo da sequência 198, 187, 176, ...

---
tipo: Discursiva
dificuldade: Média
nivel: médio
disciplina: Matemática
assunto: Análise Combinatória
gabarito:
resposta: 3840\ \text{maneiras distintas}
tags: [permutações]
fonte: concurso
concurso:
banca: UFF
ano: 1999
numero: 8
cargo:
prova:
---
\question
Cinco casais vão-se sentar em um banco de 10 lugares, de modo que cada casal permaneça sempre junto ao sentar-se.

Determine de quantas maneiras distintas todos os casais podem, ao mesmo tempo, sentar-se no banco.

---
tipo: Discursiva
dificuldade: Média
nivel: superior
disciplina: Matemática
assunto: Funções
gabarito:
resposta: I)\ \text{Verdadeira};\ II)\ \text{Falsa};\ III)\ \text{Falsa};\ IV)\ \text{Falsa};\ V)\ \text{Verdadeira}
tags: [logaritmo, exponencial, trigonometria, radiciação]
fonte: concurso
concurso:
banca: UFF
ano: 1999
numero: 9
cargo:
prova:
---
\question
Classifique cada afirmativa abaixo, em verdadeira ou falsa, justificando.

\begin{parts}
\part Para todo $x \in \mathbb{R}$, se $x < 0$, então $\sqrt{-x}$ sempre existe em $\mathbb{R}$.
\part Para todo $x \in \mathbb{R}$, $\log(-x)$ não existe em $\mathbb{R}$.
\part Para todo $x \in \mathbb{R}$, se $(x - a)^2 = (x - b)^2$ então $a = b$.
\part Para todo $x \in \mathbb{R}$, $2^{-x} < 0$.
\part Para todo $x \in \mathbb{R}$, $|\sen x| \le 1$.
\end{parts}

---
tipo: Discursiva
dificuldade: Média
nivel: superior
disciplina: Matemática
assunto: Matrizes
gabarito:
resposta: x = 0, x = -1 ou x = 1
tags: [determinante, figura]
fonte: concurso
concurso:
banca: UFF
ano: 1999
numero: 10
cargo:
prova:
---
\question
Determine o(s) valor(es) de $x$ para que a matriz

\includegraphics{figura-4}

não admita inversa.

---
tipo: Discursiva
dificuldade: Fácil
nivel: médio
disciplina: Matemática
assunto: Vetores
gabarito:
resposta: 5\ \text{m}
tags: [soma de vetores, retângulo, figura]
fonte: concurso
concurso:
banca: UFF
ano: 1999
numero: 11
cargo:
prova:
---
\question
Considere o retângulo $ABCD$ de dimensões $BC = 3$ m e $CD = 4$ m.

\includegraphics{figura-5}

Calcule $\| \overrightarrow{AB} + \overrightarrow{BD} + \overrightarrow{DC} \|$.

---
tipo: Discursiva
dificuldade: Fácil
nivel: médio
disciplina: Matemática
assunto: Geometria Analítica
gabarito:
resposta: r = 3 ou r = -3
tags: [distância entre pontos]
fonte: concurso
concurso:
banca: UFF
ano: 1999
numero: 12
cargo:
prova:
---
\question
Determine o(s) valor(es) que $r$ deve assumir para que o ponto $(r, 2)$ diste cinco unidades do ponto $(0, -2)$.

---
tipo: Discursiva
dificuldade: Média
nivel: superior
disciplina: Matemática
assunto: Polinômios
gabarito:
resposta: t = 1,\ s = 3\ \text{e}\ r = -4
tags: [resto da divisão]
fonte: concurso
concurso:
banca: UFF
ano: 1999
numero: 13
cargo:
prova:
---
\question
Determine as constantes reais $r$, $s$ e $t$ de modo que o polinômio $p(x) = rx^2 + sx + t$ satisfaça às seguintes condições:

\begin{parts}
\part $p(0) = 1$;
\part a divisão de $p(x)$ por $x^2 + 1$ tem como resto o polinômio $3x + 5$.
\end{parts}

---
tipo: Discursiva
dificuldade: Média
nivel: médio
disciplina: Matemática
assunto: Inequação Racional
gabarito:
resposta: x < -4 ou x > -2
tags: []
fonte: concurso
concurso:
banca: UFF
ano: 1999
numero: 14
cargo:
prova:
---
\question
Resolva, em $\mathbb{R} - \{-4, -2\}$, a inequação

$\frac{x-4}{x+2} < \frac{x-2}{x+4}$

---
tipo: Discursiva
dificuldade: Média
nivel: médio
disciplina: Matemática
assunto: Teorema de Pitágoras
gabarito:
resposta: AE = 4\sqrt{5}\ \text{m}
tags: [segmentos perpendiculares, figura]
fonte: concurso
concurso:
banca: UFF
ano: 1999
numero: 15
cargo:
prova:
---
\question
Na figura abaixo, os segmentos de reta $AB$, $BC$, $CD$ e $DE$ são tais que $AB$ é perpendicular a $BC$, $BC$ é perpendicular a $CD$ e $CD$ é perpendicular a $DE$.

\includegraphics{figura-6}

As medidas de $AB$, $BC$, $CD$ e $DE$ são respectivamente, 3 m, 4 m, 1 m e 4 m.

Determine a medida do segmento $AE$.

---
tipo: Discursiva
dificuldade: Média
nivel: médio
disciplina: Matemática
assunto: Geometria Plana
gabarito:
resposta: BE = 2\sqrt{21}\ \text{m}
tags: [triângulo equilátero, figura]
fonte: concurso
concurso:
banca: UFF
ano: 1999
numero: 16
cargo:
prova:
---
\question
Na figura abaixo, os triângulos $ABC$ e $DEF$ são equiláteros.

\includegraphics{figura-7}

Sabendo que $AB$, $CD$ e $DE$ medem, respectivamente, 6 m, 4 m e 4 m, calcule a medida de $BE$.

---
tipo: Discursiva
dificuldade: Difícil
nivel: superior
disciplina: Matemática
assunto: Geometria Espacial
gabarito:
resposta: h = \frac{r\sqrt{3}}{\sqrt[3]{2}}
tags: [cone, volume]
fonte: concurso
concurso:
banca: UFF
ano: 1999
numero: 17
cargo:
prova:
---
\question
Considere um cone equilátero de raio $r$ e volume $V$. Seccionou-se este cone a uma distância $h$ do seu vértice por um plano paralelo à sua base; obteve-se, assim, um novo cone de volume $\frac{V}{2}$.

Expresse $h$ em termos de $r$.

---
tipo: Discursiva
dificuldade: Difícil
nivel: superior
disciplina: Matemática
assunto: Funções
gabarito:
resposta: a)\ A(x) = -\left(\frac{\pi+4}{2}\right)x^2 + 4x;\ b)\ x = \frac{4}{\pi+4}\ \text{m}
tags: [máximos e mínimos, área, figura]
fonte: concurso
concurso:
banca: UFF
ano: 1999
numero: 18
cargo:
prova:
---
\question
Deseja-se construir uma janela com a forma de um retângulo encimado por uma semicircunferência de raio $x$ como indica a figura.

\includegraphics{figura-8}

Sabendo que o perímetro da janela deve ser igual a 4 m:
\begin{parts}
\part Expresse a área da janela em função de $x$.
\part Encontre o valor de $x$ para o qual a área da janela é a maior possível.
\end{parts}

---
tipo: Discursiva
dificuldade: Média
nivel: superior
disciplina: Matemática
assunto: Função Exponencial
gabarito:
resposta: x = 2, y = 3 ou x = 3, y = 2
tags: [sistema]
fonte: concurso
concurso:
banca: UFF
ano: 1999
numero: 19
cargo:
prova:
---
\question
Resolva o sistema

\[
\begin{cases}
3^x + 3^y = 36 \\
3^{x+y} = 243
\end{cases}
\]
